\section{Limitations}
\label{sec:limitations}

\begin{itemize}
    \item \textbf{Edits remain diffuse.} No arm we trained moved edit concentration meaningfully; every one sits in the same 0.21--0.32 band regardless of what changed underneath it. Whether diffuseness is fixable at all within a gradient-guidance approach is open.
    \item \textbf{n$\to$p stays adversarial.} On the most independent validator, both policies sit confidently below the 0.15 threshold. The ensemble-guided policy is the better of the two, but that comparison is a single run per arm (below); we cannot show the gap closes. The 0.15 threshold itself is not formally calibrated (\S\ref{sec:cv-overstates} reports counts from 0.10 to 0.25).
    \item \textbf{FOP is not computable offline}, so IP measures a related but not identical property, and SCP is a coarse proxy rather than a trained detector; in the committed configuration it is degenerate rather than merely coarse (\S\ref{sec:clinical-proxies}).
    \item \textbf{Sample size limits several verdicts}, including the 3-model ensemble's p$\to$n results ($n=32$, 27 annotated) and LHR comparison ($n=12$ shared annotated images) and FID, whose variance at $n \sim 200$ in a 1024-d feature space is high.
    \item \textbf{One training run per arm.} The single-classifier and ensemble runs share split, seeds, reward weights, action ranges and anchoring (configuration identical except the guidance flag), but differ in guidance, in classifier-selected groups (78 of 200 p$\to$n sources shared) and in hardware, and no trained-arm gap has a run-to-run error bar. In the single-classifier, 3-model and sparsity arms, the p$\to$n policy is evaluated with the n$\to$p generator adapter rather than the one it was trained with, whereas the committed arm uses its own (\S\ref{sec:directional-policies}), so their paired comparisons against the default also differ in adapter; the Grad-CAM arm's n$\to$p policy was trained for updates 1 to 49 under a different adapter from the one used at evaluation. Re-running the identical protocol on different hardware regenerates different images, so aggregate metrics move by small amounts. % TODO(noise)
Generator fine-tuning, PPO rollouts and the validity probe use only the 80\% training fold, and every evaluated image is in the 20\% test fold. The guidance classifiers and the FID extractor are pretrained models, and the \texttt{rsna} classifier was developed on RSNA data, so we cannot rule out that they have seen evaluation images.
    \item \textbf{The learned schedule is close to fixed.} The committed n$\to$p policy holds its guidance-bias control at 1.298 of 1.30 with three passes per step, and episodes last one to two steps on average, so its behaviour is close to a fixed, aggressive one-shot schedule (training-time diagnostics on 20 sampled episodes; evaluation-time action statistics are not reported). The trained-vs-random comparison ($d = 0.189$, p$\to$n only, $n = 183$) covers only the single-classifier policy, and no comparison against a tuned fixed schedule is available, so the contribution of the learned schedule itself is not isolated. The 37.9M-parameter policy network was trained on at most $800 \times 25 = 20{,}000$ transitions per direction, far fewer given the short episodes, so it may be heavily over-parameterized for this budget. % TODO(J6)
    \item \textbf{Two independent policies} rather than one target-conditioned policy; the harder direction cannot borrow capacity from the easier one.
\end{itemize}
